\chapter{Dữ liệu và thiết kế thực nghiệm}
\label{ch:experiments}

Luận văn đánh giá trên bộ dữ liệu riêng của phân xưởng, phục vụ trực tiếp bài toán ứng dụng, và bộ dữ liệu công khai AVA~v2.2. Chương này mô tả hai bộ dữ liệu, cách chia dữ liệu, cấu hình huấn luyện, giao thức đánh giá, ma trận ablation và các mô hình so sánh.

\section{Bộ dữ liệu phân xưởng}
\label{sec:factory-data}

Bộ dữ liệu được thu từ hệ thống camera giám sát của phân xưởng và gán nhãn theo cấu trúc của AVA: tại mỗi khung hình khóa, mỗi người có một hộp bao và một tập con của bốn nhãn hành vi (\cref{tab:factory-classes}); một người có thể mang nhiều nhãn đồng thời. Vì cùng cấu trúc với AVA, bộ dữ liệu dùng chung pipeline đọc dữ liệu và bộ đánh giá frame-mAP@0,5 của MMAction2. Báo cáo bổ sung \cite{baocao210726} thống kê 24.614 lượt gán nhãn, gồm 19.691 lượt cho huấn luyện và 4.923 lượt cho kiểm tra; phân bố theo lớp được trình bày trong \cref{tab:factory-stats}.

\begin{table}[H]
\centering
\caption{Bốn hành vi mục tiêu của bộ dữ liệu phân xưởng}
\label{tab:factory-classes}
\begin{tabularx}{\textwidth}{llX}
\toprule
\textbf{Nhãn} & \textbf{Thang thời gian} & \textbf{Tín hiệu nhận dạng chính}\\
\midrule
\path{sleep_on_table} & Dài hạn & Trạng thái gần như tĩnh, kéo dài nhiều giây; quan hệ đầu--thân với mặt bàn\\
\path{feet_on_table} & Dài hạn & Tư thế bền vững; bàn chân và bề mặt đỡ\\
\path{use_mobilephone} & Ngắn hạn & Vật thể nhỏ, dễ bị bàn tay che; pha đưa điện thoại lên diễn ra nhanh\\
\path{dangerous_action} & Ngắn hạn & Thay đổi nhanh trong vài khung hình; phụ thuộc thiết bị xung quanh\\
\bottomrule
\end{tabularx}
\end{table}

\begin{table}[H]
\centering
\caption{Phân bố lượt gán nhãn của bộ dữ liệu phân xưởng theo báo cáo bổ sung}
\label{tab:factory-stats}
\begin{tabularx}{\textwidth}{Xrrr}
\toprule
\textbf{Hành vi} & \textbf{Huấn luyện} & \textbf{Kiểm tra} & \textbf{Tổng}\\
\midrule
\path{use_mobilephone} & 10.477 & 2.564 & 13.041\\
\path{sleep_on_table} & 663 & 170 & 833\\
\path{feet_on_table} & 3.749 & 952 & 4.701\\
\path{dangerous_action} & 4.802 & 1.237 & 6.039\\
\midrule
\textbf{Tổng} & \textbf{19.691} & \textbf{4.923} & \textbf{24.614}\\
\bottomrule
\end{tabularx}
\end{table}

Trong tập huấn luyện, \path{use_mobilephone} chiếm khoảng 53,2\%, còn \path{sleep_on_table} chỉ chiếm 3,4\%; tỷ lệ số nhãn giữa hai lớp là khoảng 15,8 lần. Phân bố lớp giữa huấn luyện và kiểm tra tương đối ổn định, nhưng điều này không thay thế kiểm tra tính độc lập theo video hoặc camera. Các con số trên là lượt gán nhãn hành vi, không phải số khung hình khóa hay số hộp người duy nhất, vì một người có thể mang nhiều nhãn.

Báo cáo nguồn chưa thống kê số camera, video hoặc phiên ghi, tổng thời lượng, độ phân giải và tốc độ khung hình gốc, số khung hình khóa và hộp người duy nhất, cũng như định danh bộ phát hiện người. Do đó các đại lượng này không được suy ra từ số lượt gán nhãn.

\section{Cách chia dữ liệu và phòng chống rò rỉ}
\label{sec:splits}

Luận văn dùng hai cách chia bộ dữ liệu phân xưởng:
\begin{itemize}
  \item \textbf{S1 --- chia ngẫu nhiên}, như các lượt chạy ban đầu của hồ sơ thực nghiệm. Cách chia này cho phép đối chiếu với kết quả đã có, nhưng cùng một cảnh và cùng một người có thể xuất hiện ở cả tập huấn luyện lẫn tập kiểm tra, nên kết quả có xu hướng lạc quan.
  \item \textbf{S2 --- chia theo camera}: giữ riêng hai camera chỉ dùng để kiểm tra, các camera còn lại dùng để huấn luyện. S2 tránh rò rỉ cùng cảnh, cùng người và mô phỏng tình huống đưa hệ thống vào vận hành trên camera mới.
\end{itemize}
Mọi cấu hình ablation được chạy trên cả hai cách chia; chênh lệch giữa S1 và S2 ước lượng mức phụ thuộc của từng thành phần vào cảnh đã thấy khi huấn luyện (\cref{sec:cross-camera}).

Các số 19.691 và 4.923 trong \cref{tab:factory-stats} thuộc phân chia huấn luyện/kiểm tra được báo cáo trong tài liệu bổ sung; tài liệu chưa cung cấp bảng đối chiếu số khung hình khóa và số lượt gán nhãn riêng cho S1 và S2.

Về phòng chống rò rỉ dữ liệu, rò rỉ xảy ra khi cùng một cảnh, cùng một người hoặc các khung hình gần nhau của cùng một video nằm ở cả tập huấn luyện và tập kiểm tra. S2 loại bỏ rò rỉ theo cảnh bằng cách giữ riêng hai camera; mọi siêu tham số và ngưỡng theo lớp được chọn mà không dùng tập kiểm tra. Để kiểm soát chặt hơn, cần duy trì bảng kê \code{video\_id}, \code{camera\_id}, \code{session\_id}, thời gian và checksum, đồng thời dùng băm cảm nhận hoặc embedding để phát hiện đoạn gần trùng giữa các phần chia. Với AVA, phân chia train/val chính thức tách theo video, tránh sử dụng khung hình từ cùng một phim ở cả hai phía.

\section{Bộ dữ liệu AVA v2.2}
\label{sec:ava}

AVA \cite{gu2018ava} gồm 430 video phim dài 15 phút (235 huấn luyện, 64 kiểm định, 131 kiểm tra); mỗi giây là một khung hình khóa, tại đó mọi người được gán hộp bao và một hoặc nhiều nhãn trong 80 lớp hành động nguyên tử, trong đó 60 lớp được đánh giá. Giao thức chuẩn đánh giá frame-mAP@0,5 trên tập val với hộp người từ bộ đề xuất công bố của FAIR (recall 93,9\%) \cite{feichtenhofer2019slowfast,mmaction2}.

\section{Tiền xử lý và cấu hình huấn luyện}
\label{sec:training-config}

Trên bộ dữ liệu phân xưởng, mọi cấu hình trong một cặp ablation dùng chung backbone, pipeline dữ liệu, bộ tối ưu và lịch huấn luyện; mỗi cấu hình được lặp lại với ba hạt giống. Báo cáo bổ sung \cite{baocao210726} mô tả cấu hình huấn luyện chính trong \cref{tab:factory-config}. Báo cáo chưa nêu model và số lượng GPU của các lượt chạy phân xưởng, cũng chưa cung cấp cấu hình chi tiết của từng lượt chạy ablation.

\begin{table}[H]
\centering
\caption[Cấu hình huấn luyện trên dữ liệu phân xưởng]{Cấu hình huấn luyện chính trên dữ liệu phân xưởng ghi trong báo cáo bổ sung}
\label{tab:factory-config}
\small
\begin{tabularx}{\textwidth}{lX}
\toprule
\textbf{Hạng mục} & \textbf{Giá trị}\\
\midrule
Framework & MMAction2, MMDetection và MMEngine\\
Tối ưu & AdamW; learning rate $10^{-4}$; weight decay 0,05\\
Lô / epoch & 4 mẫu/GPU; 20 epoch\\
Lịch lr & Warm-up tuyến tính từ epoch 0 đến 5; cosine annealing từ epoch 5 đến 20\\
Độ chính xác và ổn định & AMP bật; dropout 0,5; cắt gradient với chuẩn tối đa 40\\
Tăng cường dữ liệu & Scale, crop, flip ngẫu nhiên; blur, perspective, affine, color jitter và random erase\\
Phần cứng & Báo cáo chưa nêu model và số lượng GPU\\
\bottomrule
\end{tabularx}
\end{table}

\section{Giao thức đánh giá}
\label{sec:protocol}

\begin{enumerate}
  \item Mỗi cặp so sánh chỉ thay đổi đúng một yếu tố; dữ liệu, cách chia, proposal, tăng cường dữ liệu, bộ tối ưu, số epoch và bộ đánh giá giữ nguyên.
  \item Mỗi cấu hình chạy ba hạt giống; báo cáo trung bình (kèm độ lệch chuẩn khi đã tổng hợp) trên cả S1 và S2.
  \item Cùng với mAP, báo cáo Precision và Recall tại ngưỡng 0,5, tốc độ xử lý và bộ nhớ GPU; thời gian suy luận đo trên cùng GPU, cùng độ phân giải và $T$.
  \item Với pipeline dựa trên proposal, chạy thêm cấu hình dùng hộp nhãn thật làm cận trên để tách lỗi định vị khỏi lỗi phân loại.
  \item So sánh với các mô hình có kiến trúc khác phản ánh hiệu năng của toàn pipeline; đóng góp của temporal neck chỉ được kết luận từ các cặp ablation trong cùng pipeline. Số liệu công bố của STAR, OV-STAD, ST-CLIP chỉ dùng để định vị khoảng hiệu năng.
\end{enumerate}

\section{Ma trận ablation}
\label{sec:ablation}

\Cref{tab:ablation} liệt kê chín cấu hình dùng để tách đóng góp của từng thành phần; cách thực hiện từng cấu hình trong mã nguồn được nêu ở \cref{sec:gate}.

\begin{table}[htbp]
\centering
\caption{Ma trận ablation (cột ``Trạng thái'' tính đến thời điểm nộp luận văn)}
\label{tab:ablation}
\small
\begin{tabularx}{\textwidth}{lXXl}
\toprule
\textbf{Mã} & \textbf{Cấu hình} & \textbf{Mục đích} & \textbf{Trạng thái}\\
\midrule
A0 & Video Swin-T tầng 4, không neck & Baseline tối thiểu (RQ1) & S1, S2\\
A1 & A0 + temporal MHSA toàn cục & So sánh toán tử tại cùng vị trí (RQ1) & S1, S2\\
A2 & A0 + SSM một chiều & Vai trò của quét một chiều (RQ2) & S1, S2\\
A3 & A0 + Bi-SSM & Nhánh dài hạn đứng riêng (RQ2) & S1, S2\\
A4 & A0 + W-MHSA cục bộ ($W = 5$) & Nhánh ngắn hạn đứng riêng (RQ2) & S1, S2\\
A5 & Bi-SSM + W-MHSA, cộng trực tiếp & Hai nhánh, không có cổng (RQ2) & S1, S2\\
A6 & Hybrid FGTN-Mamba có cổng sigmoid (tầng 4) & Mô hình đề xuất cốt lõi (RQ1--RQ2) & S1, S2\\
A7 & A6 trên tầng 3 + tầng 4 & Đóng góp của đa tỉ lệ (RQ3) & S1, S2\\
A8 & A7 + ACRN (pipeline đầy đủ) & Đóng góp của ngữ cảnh actor--cảnh (RQ4) & S1, S2\\
% Tạm bỏ A9 khỏi bản in; giữ cấu hình để có thể khôi phục.
% A9 & A8 với BCE / focal / trọng số lớp / focal + trọng số & Đóng góp của hàm mất mát (RQ4) & Chưa có\\
\bottomrule
\end{tabularx}
\end{table}

Cặp A0 $\rightarrow$ A1 trả lời câu hỏi có cần neck không, còn A1 $\rightarrow$ A6 trả lời câu hỏi một neck tiêu chuẩn đặt ở cùng vị trí đã đủ chưa. A2, A3, A4 thử từng nhánh riêng lẻ: A2 $\rightarrow$ A3 cho biết quét ngược (dùng các khung hình sau khung hình khóa) có giá trị không; A3 và A4 là hai nhánh dài hạn và ngắn hạn khi đứng một mình. A5 và A6 cùng có hai nhánh, chỉ khác ở việc cộng với tỷ lệ cố định 0,5/0,5 hay dùng cổng sigmoid, nên chênh lệch giữa chúng là đóng góp của cổng. A6 $\rightarrow$ A7 đo đóng góp của đa tỉ lệ, A7 $\rightarrow$ A8 đo đóng góp của ACRN.

\section{Các mô hình so sánh}
\label{sec:baselines}

Trên bộ dữ liệu phân xưởng, pipeline đề xuất được so sánh với năm mô hình trong hồ sơ thực nghiệm (\cref{tab:main-results}): SlowFast-R50 \cite{feichtenhofer2019slowfast}, baseline hai giai đoạn kinh điển với backbone tích chập 3D hai luồng; YOWOv2-Nano \cite{yang2023yowov2}, phiên bản nhỏ nhất của pipeline one-stage thời gian thực YOWOv2; EfficientViT3D-b1 và FastViT3D-SA12, biến thể 3D của các backbone hiệu quả EfficientViT-B1 \cite{cai2023efficientvit} và FastViT-SA12 \cite{vasu2023fastvit}; và chính pipeline đề xuất khi bỏ ACRN.

% Bỏ mục 4.8 khỏi bản in theo yêu cầu; giữ nội dung để có thể khôi phục.
% \section{Tiêu chí chấp nhận cải tiến}
% \label{sec:acceptance}
%
% Hybrid FGTN-Mamba (A6) phải cải thiện mAP ổn định qua các hạt giống so với A0, A1 và từng nhánh riêng lẻ (A2--A4) trên cả S1 và S2, trong đó mức cải thiện so với toán tử tiêu chuẩn tại cùng vị trí (A1) cần đạt tối thiểu 3 điểm mAP. Lợi ích không được chỉ đến từ tăng tham số; tốc độ triển khai phải báo cả độ trễ nhìn trước của quét hai chiều; và mọi số liệu phải tái tạo được từ cấu hình, nhật ký và checkpoint đã lưu.

\section*{Tiểu kết chương}
\addcontentsline{toc}{section}{Tiểu kết chương}

Chương này đã mô tả bộ dữ liệu phân xưởng với bốn hành vi mục tiêu, hai cách chia dữ liệu S1 (ngẫu nhiên) và S2 (theo camera), bộ dữ liệu AVA~v2.2, cấu hình huấn luyện, giao thức đánh giá, ma trận ablation chín cấu hình và năm mô hình so sánh. \Cref{ch:results} tiếp theo trình bày kết quả và thảo luận đóng góp của từng thành phần.
